% D03-001. DATA-OZZ v2.4.3, f4993e3e...; selected source claims and limits.
\section{Подготовка поднабора осциллограмм однофазных замыканий на землю}
\label{sec:ch3_earth_fault_subset}

Унификация общего набора осциллограмм создаёт основу для дальнейших
исследований, однако для анализа конкретного процесса требуется отобрать
соответствующие записи. Одной из таких задач является изучение однофазных
замыканий на землю в кабельных сетях. В совместной работе по подготовке
этого поднабора использовался исходный массив из 50~765 записей,
описанный ранее в настоящей главе.\footnote{Евдаков А.Е., Филатова Г.А.
Подготовка и анализ набора осциллограмм для исследования однофазных
замыканий на землю в кабельных сетях. Использована версия материала v2.4.3.}
Отбор и последующая ручная проверка рассматривались как самостоятельная
задача подготовки данных, предшествующая статистическому анализу.

На первом этапе сигналы приводились к относительным единицам с помощью
коэффициентов нормализации. Для поиска предполагаемых интервалов ОЗЗ
анализировались первые гармоники напряжения нулевой последовательности
и фазных напряжений. Их значения сравнивались с заданными порогами
по амплитуде и длительности. Дополнительно исключались интервалы
со сходными амплитудами напряжений трёх фаз. Такой порядок позволял
выделить кандидатов для дальнейшего просмотра, но сам по себе
не подтверждал наличие ОЗЗ в каждой выбранной записи.

Все отобранные кандидаты проверялись вручную. При просмотре исключались
записи, ошибочно принятые за ОЗЗ, уточнялись границы процесса
и оценки перенапряжений. Максимальное относительное фазное напряжение
рассматривалось в выделенной зоне ОЗЗ. Коммутационные броски при
ликвидации повреждения, находившиеся за её пределами, в эту оценку
не включались. Поэтому исследуемый максимум не следует отождествлять
с максимальным напряжением всей осциллограммы.

\begin{table}[htbp]
\centering
\caption{Распределение кандидатов и проверенного поднабора по кратности перенапряжения}
\label{tab:ch3_earth_fault_selection}
\begin{tabular}{|c|r|r|}
\hline
Кратность & До ручной проверки & После ручной проверки \\
\hline
$<1{,}2$ & 484 & 10 \\
$1{,}2\ldots1{,}71$ & 401 & 83 \\
$1{,}71\ldots1{,}75$ & 37 & 17 \\
$1{,}75\ldots2$ & 550 & 403 \\
$2\ldots2{,}5$ & 318 & 286 \\
$2{,}5\ldots3$ & 45 & 31 \\
$3\ldots3{,}5$ & 0 & 0 \\
$3{,}5\ldots4$ & 8 & 0 \\
$4\ldots4{,}5$ & 18 & 0 \\
$4{,}5\ldots5$ & 4 & 0 \\
$>5$ & 0 & 0 \\
\hline
Всего & 1865 & 830 \\
\hline
\end{tabular}
\end{table}

В таблице~\ref{tab:ch3_earth_fault_selection} подсчитаны осциллограммы,
распределённые по диапазонам максимальной кратности перенапряжения
в выбранных зонах. Сумма строк до ручной проверки составляет 1865,
после неё в поднаборе остаются 830 записей.\footnote{В тексте и итоговой
строке исходной статьи указано 1866 кандидатов. Сумма её одиннадцати
числовых строк равна 1865. Для диссертации принято значение,
подтверждённое суммированием; исходное расхождение сохранено.}
При уточнении границ и напряжений запись могла перейти в другой
диапазон. Разность численностей отдельной строки поэтому не является
самостоятельной оценкой числа ошибок алгоритма. Итоговые 830 записей
также не характеризуют полноту поиска всех ОЗЗ в исходном архиве,
поскольку весь массив не проверялся вручную на отсутствие этих процессов.

При интерпретации результатов необходимо учитывать ограничения
измерительных цепей. Нелинейность трансформаторов напряжения и
ограничение диапазона регистрации могут изменять форму сигнала
и занижать зарегистрированные пики. Неисправность измерительной цепи
может, напротив, создавать завышенные значения, не соответствующие
напряжению исследуемого сетевого процесса. Вследствие этого одного
численного максимума недостаточно для подтверждения физического
смысла записи. Ручная проверка дополняет расчёт рассмотрением
формы сигналов и последовательности событий.

Происхождение исходных записей также влияет на состав поднабора.
Осциллограммы, зарегистрированные устройствами БАВР, смещают его
в сторону повреждений, которые развиваются в многофазные короткие
замыкания. Такой материал полезен для изучения последовательности
аварийного процесса, но его распределение нельзя автоматически
распространять на все ОЗЗ в сети. Наличие отдельных развивающихся
повреждений не устанавливает вероятность их перехода в КЗ
и не подтверждает готовый метод прогнозирования.

Полученный поднабор позволяет перейти от общей подготовки данных
к исследованию зарегистрированных перенапряжений и сложностей
распознавания ОЗЗ. Эти задачи отличаются от определения направления
мощности и применимости РНМ. Для проверки соответствующих органов
необходимы собственные целевые состояния и отдельные условия оценки.
